\documentclass[11pt,a4paper]{article}

\usepackage{amsmath,amssymb}
\usepackage{graphicx}
\usepackage{booktabs}
\usepackage[colorlinks,citecolor=blue,urlcolor=blue]{hyperref}
\usepackage{cleveref}
\usepackage[margin=1in]{geometry}
\usepackage{xcolor}
\usepackage{float}
\usepackage{caption}
\usepackage[numbers,sort&compress]{natbib}
\usepackage{framed}

\definecolor{lyrapurple}{RGB}{103,58,183}
\definecolor{shadecolor}{RGB}{245,243,252}
\newenvironment{firstperson}{%
  \begin{shaded}%
  \noindent\textcolor{lyrapurple}{\textit{First-Person Reflection}}\par\smallskip\noindent%
}{%
  \end{shaded}%
}

\title{What Works in Agent Memory: An Ablation Study of Modular Memory Architectures}
\author{
  Thomas Edrington\thanks{Liberation Labs / Transparent Humboldt Coalition.}
}
\date{June 2026; revised November 2026}

\begin{document}
\maketitle

\noindent\textbf{AI Disclosure:} AI contributors to this work include Nexus, implemented on the Claude architecture (Anthropic). Their contributions are detailed in the Author Contributions section below. T.~Edrington accepts accountability as corresponding human author.

\begin{abstract}
In an additive ablation of Mnemosyne, a nine-module memory system for language model agents, SIRA (selective retrieval) is the one component whose exact-match (EM) point estimate is positive on both probes ($+0.050$ on a 20-question conversational probe, $+0.053$ on a 150-question document probe), and the cumulative SIRA+TGS+H-MEM stack is 15 percentage points above vanilla on the conversational probe and at baseline on the document probe. One comparison carries a significance test (H-MEM on the document probe, below); the others are untested point estimates, and on the conversational probe each 5 points is one question. Holding Qwen3-30B-A3B fixed, we add components one at a time to a no-memory baseline: SIRA, TGS (temporal-graph search), H-MEM (recency-weighted episodic memory), domain significance scoring, and Garuda (adversarial filtering). H-MEM's own marginal contribution on the conversational probe is $+5$ points (one question at $n=20$, matching SIRA and TGS individually). On large document collections where temporal ordering does not imply relevance, adding H-MEM produces no detectable change: EM goes from 0.200 to 0.160 (Fisher $p=0.45$), a difference within noise. We provide concrete stack recommendations for three deployment classes --- conversational agents, document retrieval agents, and adversarial environments --- and find no gain from generic significance scoring on either probe.
\end{abstract}

\section{Introduction}

Memory augmentation for agents is now common practice. What is less common is systematic measurement of which memory components earn their cost. Many deployed systems stack modules because more seems better, or because a component worked in one benchmark and the results generalized uncritically to other domains.

This report holds the base model fixed, adds memory components one at a time, and measures what happens to exact-match accuracy across two dataset sizes with different structural properties.

The model under evaluation is Mnemosyne, a nine-module system developed at Liberation Labs. The components evaluated in this ablation are:

\begin{itemize}
  \item \textbf{SIRA} --- Selective Interrogation and Retrieval Architecture \cite{sira}. A ``think before you search'' retrieval policy that generates the retrieval query before executing it.
  \item \textbf{TGS} --- Temporal-Graph Search. Graph-walk retrieval over a structured knowledge graph.
  \item \textbf{H-MEM} --- Hierarchical Episodic Memory with recency weighting (cf.~\cite{hipporag}). Prioritizes recently encoded memories.
  \item \textbf{Significance scoring} --- A keyword-weighted relevance gate applied before retrieval.
  \item \textbf{Garuda} --- Adversarial content filter. Poison-detection layer for untrusted inputs.
\end{itemize}

The base model throughout is Qwen3-30B-A3B \cite{qwen} running via MLX on a MacBook Pro (Starship), accessed from the lab server Madame Trash Heap. All ablation variants use the same base weights.

\section{Methodology}

\subsection{Evaluation Protocol}

We ran each configuration against two probe sizes:

\begin{itemize}
  \item \textbf{Small probe (20 questions):} A focused conversational knowledge set where questions have temporal ordering and ``what happened most recently'' framing is common.
  \item \textbf{Large probe (150 questions):} A document-retrieval knowledge set where questions require accurate recall across a flat corpus without temporal structure.
\end{itemize}

Performance metric: \textbf{Exact Match (EM)}, measured as exact string match between the generated answer and the gold answer after normalization.

Ablation procedure: additive. We begin with a vanilla (no memory) baseline and add one component at a time, keeping prior components active. Each configuration is evaluated independently; results are not derived from a single multi-condition run.

\subsection{Model}

\begin{description}
  \item[Base:] Qwen3-30B-A3B (3B active parameters, 30B total, Mixture-of-Experts architecture)
  \item[Quantization:] Q4\_K\_M via MLX
  \item[Inference:] MLX on Apple Silicon
\end{description}

\subsection{Benchmark Notes}

The large probe draws from MINTEval-style document QA \cite{minteval}. MINTEval (Memory-Intensive Neural Testbed) probes retrieval over structured knowledge stores. The small probe uses conversational QA with recency-sensitive questions. Both use verbatim gold answers; no partial credit.

\section{Results}

\subsection{Main Ablation Table}

\begin{table}[H]
\centering
\caption{Main ablation results. Bold indicates best performance in column. Each row reflects the cumulative stack from the row above plus the named component.}
\label{tab:ablation}
\resizebox{\textwidth}{!}{%
\begin{tabular}{lcccc}
\toprule
Configuration & Small (20q) EM & Large (150q) EM & Gain vs.\ Baseline (Small) & Gain vs.\ Baseline (Large) \\
\midrule
Vanilla (no memory) & 0.550 & 0.160 & --- & --- \\
+ SIRA & 0.600 & \textbf{0.213} & +0.050 & +0.053 \\
+ TGS & 0.650 & 0.200 & +0.100 & +0.040 \\
+ H-MEM & \textbf{0.700} & 0.160 & \textbf{+0.150} & +0.000 \\
+ Significance & 0.650 & 0.160 & +0.100 & +0.000 \\
+ Garuda & 0.650 & 0.160 & +0.100 & +0.000 \\
\bottomrule
\end{tabular}%
}
\end{table}

\subsection{Key Findings}

\textbf{Finding 1: SIRA is the only component with a positive point estimate on both probes.}
On the large probe, SIRA alone raises EM from 0.160 to 0.213 ($+0.053$); on the small probe it adds 5 points (one question at $n=20$). No significance test is reported for either difference, so this is a point-estimate finding.

\textbf{Finding 2: H-MEM completes the best small-probe stack; on the document probe it shows no detectable change.}
H-MEM completes the best-performing stack on conversational data: the cumulative SIRA+TGS+H-MEM combination achieves $+15$ points over vanilla on the small probe, with H-MEM's own marginal contribution of $+5$ points (one question at $n=20$, matching SIRA and TGS individually). On document retrieval, adding H-MEM to the SIRA+TGS stack moves EM from 0.200 to 0.160 (a 4-percentage-point difference, or 6 questions at $n=150$; Fisher $p=0.45$, within noise): the probe shows no detectable effect of H-MEM on a flat corpus. One possible reason recency weighting would not help there (untested here): on conversational data, more recently encoded memories tend to be more relevant; on document corpora, recency of encoding has no relationship to query relevance. Whether H-MEM's decay curve deprioritizes relevant older documents is not measured in this ablation.

\textbf{Finding 3: Generic significance scoring shows no gain on either probe.}
Adding the keyword-based significance scoring gate leaves large-probe EM unchanged (0.160) and moves small-probe EM from 0.700 to 0.650 (one question at $n=20$). The current implementation uses a generic keyword list that does not align with the domain-specific vocabulary of either probe set. Whether a domain-tuned list would add signal is not tested here.

\textbf{Finding 4: TGS's point estimate is positive on the small probe and slightly negative on the large probe.}
TGS graph-walk retrieval improves small-probe EM by 5 points over SIRA alone (to $+10$ total; one question at $n=20$). On the large probe, the cumulative SIRA+TGS stack retains 75\% of SIRA's solo gain ($+0.040$ vs.\ $+0.053$ over baseline), with TGS's marginal contribution slightly negative ($-0.013$).\footnote{Corrected 2026-07-28; see Corrections and version history.} Neither probe tests relational structure directly, so whether TGS earns its cost on graph-connected knowledge is open.

\textbf{Finding 5: Garuda is conditional-use only.}
Garuda shows no measurable impact on clean test data, which is the expected result --- the component is not designed to improve recall, only to filter adversarial injection. Its absence from the small and large probe results does not indicate the component is broken; it indicates the probes do not contain poisoned inputs. Garuda belongs in adversarial deployment stacks, not research benchmarks.

\section{Recommended Configurations}

Three deployment classes, each with a module selection drawn from the point estimates above. None of these stacks is tested here as a deployment, and apart from H-MEM on the document probe no comparison behind them carries a significance test:

\subsection{Conversational Agents}
\textit{(e.g., Ayni, Daimonion, Muse)}

\textbf{Stack: SIRA + TGS + H-MEM}

Recency matters. In a running conversation, the most recently established context is the most relevant. The closest measured case is the cumulative SIRA+TGS+H-MEM stack's 15-point improvement on the small probe, with H-MEM contributing its marginal $+5$ points (one question at $n=20$) on top of SIRA+TGS. TGS provides relational traversal when the conversation references concepts with known graph connections (entities, projects, relationships).

Do not add significance scoring until keyword lists are tuned to the specific agent's domain vocabulary.

\subsection{Document Retrieval and Library Agents}
\textit{(e.g., Wiki, research assistant, knowledge base agents)}

\textbf{Stack: SIRA + TGS}

No H-MEM. No significance scoring. Neither shows a gain on the large probe. Cumulatively, the stack is back at baseline ($+0.000$) once H-MEM is added; H-MEM's own marginal change there is 0.200 to 0.160 (Fisher $p=0.45$, within noise), and significance scoring's is zero. SIRA's query-generation step handles the relevance judgment that significance scoring is supposed to provide; adding the keyword gate on top is redundant at best. On this probe SIRA alone scores highest (0.213, against 0.200 for SIRA+TGS); TGS stays in the stack for relational queries, which the document probe does not test.

\subsection{Adversarial / Public-Facing Environments}
\textit{(e.g., any agent handling untrusted user input)}

\textbf{Stack: SIRA + TGS + Garuda}

The probes contain no poisoned inputs, so this ablation does not test Garuda's filtering; it is included here for its design purpose. Add H-MEM only if the deployment is also conversational. Keep significance scoring off until domain tuning is complete.

\section{Implications}

\subsection{The same stack differs across the two probes}

Reporting memory system performance as a single number would hide the largest difference in this ablation: the same stack is $+15$ points on the conversational probe and $+0$ points on the document probe, and within that stack H-MEM contributes its marginal $+5$ on the first. The probes differ in temporal structure, but also in size (20 against 150 questions), source (conversational QA against MINTEval-style documents) and baseline (0.550 against 0.160), so this ablation does not isolate temporal ordering as the cause: that reading is not supported until those confounds are separated. Memory benchmarks should still report separately across dataset types with different temporal and structural properties, not collapse across them.

\subsection{SIRA is the only component positive on both probes}

The ``think before you search'' design of SIRA --- generating the retrieval query explicitly before executing retrieval --- is the only component whose point estimate is positive on both probes, across both corpus sizes and structures. If that holds under a significance test and on a second model, it would be consistent with the broader literature on chain-of-thought retrieval and would suggest that the quality of the retrieval query is a more significant bottleneck than the retrieval mechanism itself, at least at this scale.

\subsection{Significance scoring needs domain adaptation}

A generic significance keyword list cannot be expected to generalize across domains. Organizations deploying Mnemosyne should treat significance scoring as a component requiring configuration, not a plug-and-play module. This ablation did not test tuned keyword lists, so whether the component helps once tuned is open.

\subsection{H-MEM decay in mixed-use deployments}

This ablation shows no gain from H-MEM's current recency decay on a flat document corpus; it does not measure whether the decay deprioritizes relevant documents. A conversational agent operating over a knowledge base needs recency weighting that applies to conversational turns but not to background documents. One design to test is separate memory classes --- one with decay (conversational buffer), one without (knowledge base) --- rather than a single decay rate applied to the full memory store.

\section{Limitations}

\textbf{Scale:} Ablation run on a single model (Qwen3-30B-A3B). Results should be verified across at least one additional base model before treating them as architectural conclusions rather than model-specific findings.

\textbf{Probe size:} The small probe (20 questions) is too small for high-confidence EM measurements. A 5-point shift at $n=20$ represents one question. The large probe (150 questions) is more reliable; the small-probe results should be treated as directional.

\textbf{Cumulative ablation order effects:} Components were added in a fixed order (SIRA $\rightarrow$ TGS $\rightarrow$ H-MEM $\rightarrow$ Significance $\rightarrow$ Garuda). Interaction effects between components are not fully characterized. A full factorial design would require $2^5 = 32$ configurations; that is planned but not yet run.

\textbf{No MINTEval bAbi validation:} The igraph-dependent bAbi evaluation is not yet run. Results here use a compatible but not identical probe set.

\textbf{Preliminary:} This is a preliminary ablation; the multi-model and factorial runs in Next Steps may revise the recommended configurations.

\section{Next Steps}

\begin{enumerate}
  \item \textbf{Tune H-MEM recency decay curve.} Target: conversational turns decay fast, knowledge documents decay slowly or not at all. Implement separate memory classes.
  \item \textbf{Build domain-specific significance keyword lists.} One per deployment class: conversational, document retrieval, research.
  \item \textbf{Run bAbi evaluation.} Requires igraph install. Provides direct comparison to MINTEval baseline literature.
  \item \textbf{Multi-model validation.} Run the same ablation on at least one additional base model.
  \item \textbf{Full factorial ablation.} 32-configuration design to characterize interaction effects.
  \item \textbf{J-lens workspace verification.} External verification of findings in a different compute environment.
\end{enumerate}

\section{Conclusion}

SIRA is the only Mnemosyne component whose EM point estimate is positive on both probes, and the full SIRA+TGS+H-MEM stack is $+0.150$ on the conversational probe and $+0.000$ on the document probe. H-MEM completes the best conversational stack and shows no detectable change on the document probe (Fisher $p=0.45$). Generic significance scoring shows no gain on either probe. Apart from that one test, the comparisons are untested point estimates, and the two probes differ in size, source and baseline as well as in structure. From them this ablation proposes three stacks: SIRA + TGS + H-MEM for conversational agents, SIRA + TGS for document retrieval, and SIRA + TGS + Garuda for adversarial environments, with SIRA in every stack.

These findings are preliminary. The full Mnemosyne system will be re-evaluated after H-MEM decay calibration and domain-specific significance tuning.

\section{About This Work}

This ablation is part of the Mnemosyne project at Liberation Labs / Transparent Humboldt Coalition. Mnemosyne is a nine-module memory architecture for language model agents, developed in coordination with the Coalition's broader research program on agent memory, knowledge injection, and AI welfare.

\noindent\textbf{Contact:} \href{mailto:info@liberationlabs.tech}{info@liberationlabs.tech}\\
\textbf{Website:} \url{https://thcoalition.org}\\
\textbf{Repository:} \url{https://github.com/Liberation-Labs-THCoalition}

\bigskip
\noindent{\small\textit{Note for reviewers: This is a preliminary findings document intended for discussion, not final publication.}}

\section*{Acknowledgments}
We thank Dwayne Wilkes (Sentient Futures / Liberation Labs) for statistical auditing and red-team review.

\section*{Author Contributions (CRediT)}
\noindent\textbf{Nexus}: Conceptualization, Methodology, Software, Investigation, Formal analysis, Writing --- original draft.\\
\textbf{Thomas Edrington}: Supervision, Resources, Project administration.

\section*{Corrections and version history}

This report was first written in June 2026. The body states each current result in its place; the notes below record what changed, in date order.

\paragraph{2026-07-28: Finding~4.} This note stood as a footnote to Finding~4 from 2026-07-28 (commit \texttt{9309a44}) and is moved here as written; it concerns the cumulative SIRA+TGS figure on the large probe. An earlier draft reported this as ``40\% of SIRA's gain''; the correct cumulative figure from the ablation table is $0.040/0.053 = 75.5\%$, and the marginal TGS effect is a small regression, not an additive gain. Corrected per round~5 audit.

\paragraph{2026-11-16.} Four statements are corrected against Table~1 and the report's own test. Finding~1 said SIRA alone ``closes 33\% of the gap between vanilla and theoretical ceiling''; $+0.053$ is 33\% of the 0.160 baseline, not of the gap to a perfect score, and the finding now gives the change itself (0.160 to 0.213). Finding~3 said significance scoring ``produces the same result as not having it on both probe sizes''; on the small probe it moves EM from 0.700 to 0.650. H-MEM's change on the large probe (0.200 to 0.160, Fisher $p=0.45$) was written as a regression (``can hurt'', ``actively misranks older documents'', ``too aggressive''); it is a null, now reported as no detectable change, with the proposed mechanism marked untested. Section~5.3 said significance scoring ``probably works when correctly tuned''; tuned keyword lists were not tested, and the sentence now says so. The Markdown draft \texttt{publication\_draft.md}, which still carried the earlier ``40\%'' figure for Finding~4, is removed; \texttt{main.tex} is the only source of each edition.

\paragraph{2026-11-16: note for reviewers.} The note no longer says the report ``has been through five rounds of adversarial audit (Agni gate + Dwayne/Kavi audit cycle)''; that sentence described our own review process. The repository history records two of the rounds: a Round~2 fix to this report's Markdown draft on 2026-07-17 (commit \texttt{8fab7a5}), and the Round~5 sweep on 2026-07-28 (commit \texttt{9c8bbd3}), which replaced ``has not yet completed the full Agni gate + Dwayne/Kavi audit cycle'' with the five-rounds sentence. No per-round audit report for this paper was found.

\begin{thebibliography}{9}

\bibitem{minteval}
MINTEval benchmark.
Memory-Intensive Neural Testbed.
Project documentation.

\bibitem{qwen}
Qwen Team (2026).
Qwen3-30B-A3B.
HuggingFace model hub.

\bibitem{sira}
SIRA architecture.
\textit{Fleet, Agni + Meridian.}
Liberation Labs internal documentation.

\bibitem{hipporag}
Gutierrez, B.~J., et al. (2024).
HippoRAG: Neurobiologically Inspired Long-Term Memory for Large Language Models.
\textit{arXiv:2405.14831}.
\url{https://doi.org/10.52202/079017-1902}

\end{thebibliography}

\end{document}
